In this subsection, we express each side of the equivalence (\ref{eqn:cta}) of \Cref{thm:Cta} as the category of modules over some commutative algebra in $\Sp_{C_2, i2} ^{\Fil}$.
Almost everything in this subsection is a standard application of Lurie's theory of higher algebra.
We have extracted the necessary material in \Cref{app:boilerplate}, where we present it in a digested form. 

We begin with the left-hand-side of (\ref{eqn:cta}).

\begin{lem} \label{lem:gr-of-comparison}
  The cofiber of $\ta : \Ss_2^{0,0,-1} \to \Ss_2^{0,0,0}$ is a commutative algebra in $\SHR_{i2}$. Furthermore, there is an equivalence of symmetric monoidal categories
  \[ \mathrm{Mod}(\SHRat; C\ta) \simeq \mathrm{Mod}( \Sp_{C_2, i2}^{\Gr}; \Gr( R_\bullet ) ), \]
	where $R_{\bullet} =\Tot^{*}\left( P_{2\bullet} \MU_{\mathbb{R},2}^{\otimes *+1} \right)$ 
\end{lem}

\begin{proof}
  Recall that the shift map $\Ss_2(-1) \to \Ss_2$ in $\Sp_{C_2, i2}^{\Fil}$ is denoted by $\tau$.
    Tensoring the equivalence of \Cref{exm:fil-ctau} with $\Sp_{C_2, i2}$ and applying \Cref{lem:tensor-mod}, we see that there is an equivalence of presentably symmetric monoidal categories
    \[ \Sp_{C_2, i2} ^{\Gr} \simeq \Mod ( \Sp_{C_2, i2} ^{\Fil} ; C\tau ). \]

    Under the symmetric monoidal functor $i^* : \Sp_{C_2, i2} ^{\Fil} \to \SHRat$ constructed in \Cref{prop:top-diagram}, the commutative algebra $C\tau$ maps to $C\ta$. Therefore, using \Cref{lem:tensor-mod}, we have 
    \[ \mathrm{Mod}(\SHRat; C\ta) \simeq \Sp_{C_2, i2}^{\Gr} \otimes_{\Sp_{C_2, i2} ^{\Fil}} \SHRat. \]
    Using \Cref{thm:filt-model}, there is an equivalence of presentably symmetric monoidal categories under $\Sp_{C_2, i2}^{\Fil}$ between $\SHRat$ and $\mathrm{Mod}( \Sp^{\Fil}_{C_2,i2} ; R_\bullet )$. Tensoring down along the associated graded ring map and using \Cref{lem:tensor-mod}, we obtain equivalences
    \begin{align*}
      \Sp_{C_2, i2}^{\Gr} \otimes_{\Sp_{C_2, i2}^{\Fil}} \SHRat
      %&\cong \Sp^{\Gr} \otimes_{\Sp^{\Fil}} \mathrm{Mod}( \Sp_{C_2,i2}^{\Fil} ; \Gamma_* \Ss_2 ) \\
      &\simeq \Sp_{C_2,i2}^{\Gr} \otimes_{\Sp_{C_2,i2}^{\Fil}} \mathrm{Mod}( \Sp_{C_2,i2}^{\Fil} ; R_\bullet ) \\
      &\simeq \mathrm{Mod}( \Sp_{C_2,i2}^{\Fil} ; \Gr(R_{\bullet}) ). \qedhere
    \end{align*}
      %
  %% Note that the shift map $\tau: \o(1) \to \o$ in $\Fil(\Sp_{C_2, i2})$ goes to $\ta$ under the colimit-preserving symmetric monoidal functor $c_{\C/\R} : \Fil(\Sp_{C_2, i2}) \to \SH(\R)^{\at} _{i2}$. Since $C\tau$ is an $\E_{\infty}$-ring in $\Fil(\Sp_{C_2, 2})$, it follows that its image $C\ta$ under $c_{\C/\R}$ is as well.

  
  %% %% \[ \mathrm{Mod}(\SHgc_{\mathrm{i}p}; C\ta) \simeq \Gr(\Sp) \otimes_{\Fil(\Sp)} \SHgc_{\mathrm{i}p} \simeq \Gr(\Sp) \otimes_{\Fil(\Sp)} \mathrm{Mod}(\Fil(\Sp_{C_2}); R_\bullet) \simeq \mathrm{Mod}(\Gr(\Sp_{C_2}); R_\bullet) \]
  %%   \[ \Gr(\Sp_{C_2, i2}) \otimes_{\Fil(\Sp_{C_2, i2})} \SH(\R)^{\at}_{i2} \simeq \Gr(\Sp_{C_2, i2}) \otimes_{\Fil(\Sp_{C_2, i2})} \mathrm{Mod}(\Fil(\Sp_{C_2,2}); \Gamma_* (\Ss_2)) \]
  %%   symmetric monoidal categories. Using \Cref{lem:tensor-mod}, we can simplify each side of this equivalence to obtain:
  %%   \[ \Gr(\Sp_{C_2, i2}) \otimes_{\Fil(\Sp_{C_2, i2})} \SH(\R)^{\at}_{i2} \simeq \mathrm{Mod}(\SH(\R)^{\at}_{i2}; C\ta), \]
  %%   \[ \Gr(\Sp_{C_2, i2}) \otimes_{\Fil(\Sp_{C_2, i2})} \mathrm{Mod}(\Fil(\Sp_{C_2}); (\Gamma_* (\Ss_2)) \simeq \mathrm{Mod}(\Gr(\Sp_{C_2}); \Gr(\Gamma_* (\Ss_2))), \]
  %%   where in the first case we have used the fact that $\Mod(\Fil(\Sp_{C_2, i2}); C\tau) \simeq \Gr(\Sp_{C_2,i2})$.
\end{proof}

Now we proceed to the right-hand-side of (\ref{eqn:cta}).

\begin{lem} \label{lem:omega-euler-one}
  The Euler characteristic of $\omega_{\G/\Mfg} \in \IndCoh(\Mfg)^{\heartsuit}$ is equal to 1.
  %% one: $\chi(\omega_{\G/\Mfg}) = 1$.
  As a consequence, \Cref{cnstr:twist} provides a symmetric monoidal left adjoint
  \[ p^* : \Ab^{\Gr, \heartsuit} \to \IndCoh(\Mfg)^{\heartsuit} \]
  %% \subset \IndCoh(\Mfg)
  which sends $\Z(1)$ to $\omega_{\G/\Mfg}$.
\end{lem}

\begin{proof}
  Let $L$ denote the Lazard ring. Then the flat cover $\Spec(L) \to \Mfg$ determines a pullback map $\IndCoh(\Mfg) \to \Mod_{L}$. This determines an injective pullback map
  $\pi_0 \End(\mathcal{O}_{\Mfg}) \to \pi_0 \End(L)$.
  Since the pullback of $\omega_{\G/\Mfg}$ is equivalent to $L$ we learn that $\chi(\omega_{\G/\Mfg}) = 1$.
  %% The latter group is isomorphic to $\Z ^\times$, which the former group contains, so the map is an equivalence.
  %% It follows that $\chi(\omega_{\G/\Mfg})$ is equal to the Euler characteristic of the pullback of $\omega_{\G/\Mfg}$ to $\Mod_{L}$. But this pullback is just $L$ itself, which implies that its Euler characteristic must be $1$, as desired.
  %% For the second statement, it is clear that one may construct a monoidal functor $\Z \to \IndCoh(\Mfg)^{\heartsuit}$ sending $1$ to $\omega_{\G/\Mfg}$. To check that it is symmetric monoidal, one has to check that it sends the swap automorphism of $2 = 1+1$, which is the identity, to the swap isomorphism of $\omega_{\G/\Mfg} ^2 \simeq \omega_{\G/\Mfg} \otimes \omega_{\G/\Mfg}$. This is precisely the condition that $\chi(\omega_{\G/\Mfg}) = 1$.
\end{proof}

Since $p^*$ sends a family of compact dualizable objects (the powers of $\Z(1)$) to a family of compact dualizable generators (the powers of $\omega_{\G/\Mfg}$), we may apply \Cref{prop:rigid} to obtain the following lemma:

\begin{lem}\label{lem:MU-koszul}
  There is an equivalence of presentably symmetric monoidal categories
  \[ \mathrm{IndCoh}(\mathcal{M}_{\mathrm{fg}}) \simeq \mathrm{Mod}(\Ab^{\Gr}; p_*\mathcal{O}_{\Mfg} ). \]
\end{lem}

%% \begin{proof}
%%   By \Cref{lem:omega-euler-one}, there is a symmetric monoidal functor $\Z \to \IndCoh(\Mfg) ^{\heartsuit}$ sending $1 \in \Z$ to $\omega_{\G/\Mfg}$.
  
%%   As a consequence, we obtain a symmetric monoidal left adjoint
%%   \[\Ab^{\Gr} \to \mathrm{IndCoh}(\mathcal{M}_{\mathrm{fg}}) \]
%%   which sends $\Z(n)$ to $\omega^{n} _{\mathbb{G}/\Mfg}$.
%%   This functor sends a family of compact dualizable generators to a family of compact dualizable generators,
%%   so by [ref] \todo{Find or write ref!} it factors through the desired symmetric monoidal equivalence
%%   \[  \mathrm{Mod}(\Ab^{\Gr}; \ExtMUt ) \simeq \mathrm{IndCoh}(\mathcal{M}_{\mathrm{fg}}). \qedhere \]
%%   %% \[ \mathrm{Mod}(\Ab^{\Gr}; \ExtMU) \to \mathrm{IndCoh}(\mathcal{M}_{\mathrm{fg}}). \]
%% \end{proof}

Using \Cref{lem:MU-koszul} and \Cref{lem:tensor-mod}, we obtain the following corollary:

\begin{cor}\label{lem:uMU-koszul}
  There is an equivalence of presentably symmetric monoidal categories
  \[ \uAb_{i2} \otimes_{\Ab} \mathrm{IndCoh}(\mathcal{M}_{\mathrm{fg}}) \simeq \Mod ( \uAb_{i2}^{\Gr}; \underline{p_*\mathcal{O}_{\Mfg}} \otimes_{\uZ} \uZt ).  \]
\end{cor}

%% \begin{proof}
%%   Using \Cref{lem:MU-koszul}, we have equivalences
%%   \begin{align*}
%%     \uAb \otimes_{\Ab} \mathrm{IndCoh}(\mathcal{M}_{\mathrm{fg}})
%%     &\simeq \uAb \otimes_{\Ab} \mathrm{Mod}(\Ab^{\Gr}; \ExtMUt ) \\
%%     &\simeq \uAb^{\Gr} \otimes_{\Ab^{\Gr}} \mathrm{Mod}(\Ab^{\Gr}; \ExtMUt ).
%%   \end{align*}
%%   We conclude by applying \Cref{lem:tensor-mod}.
%% \end{proof}

